\documentclass[aspectratio=169, 11pt, xcolor={table}]{beamer}
% Shared style for the course slides, based on presentation_text/main.tex.
% Dark blue removed: every slide (cover included) uses the light background
% of the template's content slides, with teal + orange as accents.

\usepackage{amsmath}
\usepackage{amssymb}
\usepackage{booktabs}
\usepackage{colortbl}
\usepackage{tabularx}
\usepackage{tikz}
\usepackage[skins]{tcolorbox}
\usepackage{fontawesome5}
\usepackage{polyglossia}

\usetikzlibrary{positioning, calc, arrows.meta, shapes.geometric}

%------------------------------------------------
% Languages and fonts (compile from slides/ with XeLaTeX or tectonic)
%------------------------------------------------
\setdefaultlanguage[numerals=maghrib]{arabic}
\setotherlanguage{english}
\usefonttheme{serif}

\newfontfamily\arabicfont[
    Script=Arabic, Scale=0.98, Path=../fonts/, Extension=.ttf,
    UprightFont=*-Regular, BoldFont=*-Bold,
    ItalicFont=*-Italic, SlantedFont=*-Italic,
    BoldItalicFont=*-BoldItalic, BoldSlantedFont=*-BoldItalic]{Amiri}
\newfontfamily\arabicfontsf[
    Script=Arabic, Scale=0.98, Path=../fonts/, Extension=.ttf,
    UprightFont=*-Regular, BoldFont=*-Bold,
    ItalicFont=*-Italic, SlantedFont=*-Italic,
    BoldItalicFont=*-BoldItalic, BoldSlantedFont=*-BoldItalic]{Amiri}
% Heading font
\newfontfamily\khatt[Script=Arabic, Scale=1.05, Path=../fonts/,
    BoldFont=ae_AlBattar.ttf]{ae_AlBattar.ttf}

%------------------------------------------------
% Colors (template palette without the navy tones)
%------------------------------------------------
\definecolor{brandTeal}{HTML}{00B4D8}
\definecolor{brandOrange}{HTML}{FF7043}
\definecolor{brandLight}{HTML}{F8F9FA}
\definecolor{brandDark}{HTML}{263238}
\definecolor{brandMuted}{HTML}{6C757D}
\colorlet{tealDeep}{brandTeal!70!black}

\setbeamercolor{background canvas}{bg=brandLight}
\setbeamercolor{normal text}{fg=brandDark}
\setbeamercolor{structure}{fg=brandTeal}
\setbeamertemplate{navigation symbols}{}

% beamer.cls sets \raggedright globally; under RTL that centers short lines
% or pushes them left, so go back to justified text (flush right).
\rightskip=0pt\relax
% Lists call \raggedright too; make it a no-op for RTL slides.
\AtBeginDocument{\renewcommand{\raggedright}{\rightskip=0pt\relax}}

%------------------------------------------------
% Frame title, footer, lists
%------------------------------------------------
\setbeamertemplate{frametitle}{%
    \vspace{0.35cm}
    \noindent{\khatt\bfseries\Large\color{brandDark}\insertframetitle}
    \par
    \vskip-2pt
    \noindent{\color{brandTeal}\rule{\textwidth}{1.5pt}}
    \vspace{0.05cm}
}

% Footer: page number only
\setbeamertemplate{footline}{%
    \hbox to \paperwidth{\hfill
        {\color{brandMuted}\scriptsize\LR{\insertframenumber\,/\,\inserttotalframenumber}}%
        \hspace{0.6cm}}%
    \vspace{0.25cm}}

\setbeamertemplate{itemize item}{\color{brandTeal}\small\faChevronLeft}
\setbeamertemplate{itemize subitem}{\color{brandMuted}\small\faAngleLeft}
\setbeamertemplate{enumerate item}{\color{tealDeep}\bfseries\insertenumlabel)}

% Tables: every cell framed by a visible border
\newcolumntype{R}{>{\raggedleft\arraybackslash}X}
\newcolumntype{C}{>{\centering\arraybackslash}X}
\newcommand{\ras}[1]{\cellcolor{brandTeal!20}\textbf{#1}}
% Colored text inside a table cell (\leavevmode keeps it on the row baseline)
\newcommand{\cell}[2]{\leavevmode\textcolor{#1}{#2}}
\renewcommand{\arraystretch}{1.35}
\setlength{\arrayrulewidth}{0.6pt}
\arrayrulecolor{brandMuted!70}

%------------------------------------------------
% Boxes (the two boxes of the template's 2nd slide)
%------------------------------------------------
\tcbset{
    brandbox/.style={
        enhanced, arc=4pt, outer arc=4pt, boxrule=1.5pt,
        coltitle=white, coltext=brandDark, fonttitle=\khatt\bfseries,
        toptitle=0.8mm, bottomtitle=0.8mm,
        top=1.2mm, bottom=1.2mm, left=3mm, right=3mm,
        before skip=4pt, after skip=4pt}
}
% Teal box: definitions, objectives, summaries
\newtcolorbox{tbox}[2][]{brandbox, colback=brandTeal!10,
    colframe=brandTeal, title={#2}, #1}
% Orange box: warnings
\newtcolorbox{obox}[2][]{brandbox, colback=brandOrange!10,
    colframe=brandOrange, title={#2}, #1}
% Card (template's 3rd slide, in teal)
\newtcolorbox{card}[2][]{enhanced, colback=brandTeal!6, colframe=brandTeal,
    arc=6pt, boxrule=1.2pt, halign=center, halign title=center,
    coltitle=white, coltext=brandDark, fonttitle=\khatt\bfseries,
    top=1.5mm, bottom=1.5mm, left=2mm, right=2mm, title={#2}, #1}
% Prompt text written to the tool
\newtcolorbox{amr}[1][brandTeal]{enhanced, arc=1mm, boxrule=0pt,
    colback=white, overlay={\fill[#1] (frame.north east) rectangle ([xshift=-2.5pt]frame.south east);},
    left=6pt, right=8pt, top=2pt, bottom=2pt,
    before skip=3pt, after skip=3pt}

%------------------------------------------------
% TikZ (Arabic right-to-left inside nodes)
%------------------------------------------------
\newcommand{\nodecolor}{}
\tikzset{
    every node/.append style={
        execute at begin node={\nodecolor\everypar{\beginR}\beginR},
        execute at end node=\endR},
    lawn/.code={\def\nodecolor{\color{#1}}},
    khatwa/.style={draw=#1, fill=#1!12, very thick, rounded corners=5pt,
        minimum height=0.95cm, text width=1.9cm, align=center,
        font=\khatt\bfseries\small},
    khatwa/.default=brandTeal,
    sahm/.style={-{Stealth[length=2.5mm]}, thick, color=brandMuted},
    rang/.style={draw=#1, fill=#1!12},
    raqm/.style={circle, fill=#1, lawn=white, inner sep=1.5pt,
        font=\small\bfseries},
}

% Horizontal chain of steps, right to left:
% \flow[width]{{a}, {b\\c}, {d}}
\newcommand{\flow}[2][1.9cm]{%
    \begin{tikzpicture}[node distance=5mm]
        \foreach \t [count=\i] in {#2} {
            \ifnum\i=1
                \node[khatwa, text width=#1] (f\i) {\t};
            \else
                \pgfmathtruncatemacro\p{\i-1}
                \node[khatwa, text width=#1, left=of f\p] (f\i) {\t};
                \draw[sahm] (f\p) -- (f\i);
            \fi
        }
    \end{tikzpicture}}

%------------------------------------------------
% Cover (template's 1st slide on the light background)
%------------------------------------------------
\newcommand{\cover}[3]{%
\begin{frame}[plain]
    \begin{tikzpicture}[remember picture, overlay]
        \draw[brandTeal, line width=1.5pt, opacity=0.3]
            (current page.south east) ++(-1,1) circle (4cm);
        \draw[brandTeal, line width=0.8pt, opacity=0.25]
            (current page.north west) ++(2,2) circle (6cm);
        \draw[brandOrange, line width=0.5pt, opacity=0.25]
            (current page.north west) ++(2,2) circle (7cm);
        \node[anchor=north east] at ($(current page.north east)+(-1.0,-0.8)$) {
            \begin{tikzpicture}[scale=0.35]
                \fill[brandTeal] (0,0) -- (0.8,1.2) -- (1.6,0) -- cycle;
                \fill[brandOrange] (0.8,1.2) -- (1.6,2.4) -- (2.4,1.2) -- cycle;
            \end{tikzpicture}
        };
    \end{tikzpicture}

    \vspace{-0.05cm}
    {\hspace*{1.1cm}\khatt\bfseries\color{brandDark}
        المدرسة العليا للأساتذة ببوزريعة ــ ملحقة بومرداس\par}

    \vspace{1.2cm}
    \begin{minipage}{0.85\textwidth}
        {\color{tealDeep}\khatt\bfseries #1} \\[0.3cm]
        {\raggedleft\color{brandDark}\Huge\khatt\bfseries #2\par} \vspace{0.3cm}
        {\color{brandMuted}\large #3}
    \end{minipage}

    \vspace{1.2cm}
    \begin{minipage}{0.85\textwidth}
        \color{brandMuted}\small
        \begin{tabular}{@{}ll}
            \textbf{الأستاذ:} & د. بن عبد الرزاق يوسف \\
            \textbf{الفئة:}   & طلبة قسم اللغة العربية \\
        \end{tabular}
    \end{minipage}
\end{frame}}

\title{البحث عن المعلومات، تحليلها وتلخيصها}

\begin{document}

\cover{درس الفصل الخامس}
{البحث عن المعلومات، تحليلها وتلخيصها}
{من السؤال إلى المصدر الموثوق}

\begin{frame}{أهداف الدرس}
    \begin{tbox}{\faBullseye\quad أهداف الدرس}
        في نهاية هذا الدرس، تستطيع أن:
        \begin{itemize}
            \item تفرّق بين موضوع البحث وسؤال البحث.
            \item تختار كلمات مفتاحية مناسبة.
            \item تميّز بين أنواع المصادر وتقيّم جودتها.
            \item تلخّص المعلومات بطريقة منظمة وأمينة.
        \end{itemize}
    \end{tbox}
\end{frame}

\begin{frame}{مراحل البحث عن المعلومة}
    البحث الجيد لا يبدأ بمحرك البحث، بل يبدأ \textbf{بسؤال واضح}.
    نمرّ بأربع مراحل:
    \begin{center}
        \flow[2.4cm]{{حدد\\الحاجة}, {ابحث}, {قيّم\\المصادر}, {لخّص}}
    \end{center}
\end{frame}

\begin{frame}{من الموضوع إلى سؤال البحث}
    \begin{tbox}{\faLightbulb\quad تعريف}
        \textbf{موضوع البحث:} مجال عام وواسع، مثل: «الذكاء الاصطناعي والتعليم».\\[2pt]
        \textbf{سؤال البحث:} سؤال محدد نريد الإجابة عنه، مثل:
        «كيف يساعد الذكاء الاصطناعي التلميذ على تحسين التعبير الكتابي؟».
    \end{tbox}
    \begin{center}
    \begin{tikzpicture}[scale=0.7]
        \fill[brandTeal!15] (-6,0) -- (6,0) -- (4.5,-1.1) -- (-4.5,-1.1) -- cycle;
        \fill[brandTeal!40] (-4.5,-1.1) -- (4.5,-1.1) -- (3,-2.2) -- (-3,-2.2) -- cycle;
        \fill[brandTeal]    (-3,-2.2) -- (3,-2.2) -- (1.5,-3.3) -- (-1.5,-3.3) -- cycle;
        \node[font=\khatt\bfseries\small] at (0,-0.55) {الذكاء الاصطناعي (عام جدًا)};
        \node[font=\khatt\bfseries\small] at (0,-1.65) {في التعليم الجامعي (أدق)};
        \node[font=\khatt\bfseries\small, lawn=white] at (0,-2.75) {فوائده ومخاطره};
    \end{tikzpicture}
    \end{center}
\end{frame}

\begin{frame}{مثال من اختصاصك}
    الموضوع العام يعطي نتائج كثيرة ومتفرقة، أما السؤال المحدد فيوجّه البحث.
    \begin{tbox}{\faChalkboardTeacher\quad موضوع: تعليم النحو}
        \begin{enumerate}
            \item ما الصعوبات التي يواجهها تلاميذ المتوسط في الإعراب؟
            \item هل تساعد الأمثلة من الحياة اليومية على فهم القاعدة؟
            \item كيف يمكن استعمال الذكاء الاصطناعي في تمارين الإعراب؟
        \end{enumerate}
    \end{tbox}
\end{frame}

\begin{frame}{الكلمات المفتاحية}
    \begin{tbox}{\faLightbulb\quad تعريف}
        \textbf{الكلمات المفتاحية:} الكلمات الأساسية في سؤال البحث،
        نكتبها في محرك البحث للوصول إلى المعلومة.
    \end{tbox}
    \textbf{\color{tealDeep}كيف نختارها؟}
    \begin{enumerate}
        \item نقرأ سؤال البحث، ونحذف الكلمات العامة مثل: «كيف»، «ما»، «هل».
        \item نحتفظ بالكلمات التي تحمل المعنى.
        \item نضيف المرادفات، والكلمات الإنجليزية المقابلة.
    \end{enumerate}
\end{frame}

\begin{frame}{مثال على الكلمات المفتاحية}
    \textbf{\color{tealDeep}السؤال:}
    \begin{amr}[brandTeal]
        كيف يساعد الذكاء الاصطناعي على تحسين التعبير الكتابي؟
    \end{amr}
    \textbf{\color{tealDeep}الكلمات المفتاحية:}
    \begin{center}
        \begin{tabular}{|c|c|c|c|c|}
            \hline
            الذكاء الاصطناعي & التعبير الكتابي & الإنشاء &
            \textenglish{AI} & \textenglish{Writing skills} \\ \hline
        \end{tabular}
    \end{center}
\end{frame}

\begin{frame}{أنواع المصادر ودرجة الثقة}
    ليست كل المصادر في درجة الثقة نفسها. نرتبها من الأكثر موثوقية إلى الأقل:
    \begin{center}
    \begin{tikzpicture}[
        masdar/.style={thick, rounded corners=3pt,
            minimum height=0.68cm, text width=7cm, align=center}]
        \node[masdar, rang=brandTeal]                at (0,0)     {مصدر أكاديمي: كتاب محقق، مقال علمي محكّم};
        \node[masdar, rang=brandTeal!70!brandOrange] at (0,-0.85) {موقع مؤسسة رسمية: وزارة، جامعة};
        \node[masdar, rang=brandMuted]               at (0,-1.7)  {مقال صحفي};
        \node[masdar, rang=brandOrange!70!brandTeal] at (0,-2.55) {مدونة أو موقع شخصي};
        \node[masdar, rang=brandOrange]              at (0,-3.4)  {محتوى أنتجه الذكاء الاصطناعي};
        \draw[{Stealth[length=3mm]}-, line width=2.5pt, brandMuted]
            (-4.2,0.2) -- (-4.2,-3.6);
        \node[font=\khatt\bfseries, lawn=tealDeep, anchor=east] at (-4.4,0) {ثقة أكبر};
        \node[font=\khatt\bfseries, lawn=brandOrange, anchor=east] at (-4.4,-3.4) {ثقة أقل};
    \end{tikzpicture}
    \end{center}
\end{frame}

\begin{frame}{تنبيهات حول المصادر}
    \begin{obox}{\faExclamationTriangle\quad محتوى الذكاء الاصطناعي ليس مصدرًا علميًا}
        نستعمله للفهم أو لاقتراح أفكار، ثم نبحث عن المعلومة في مصدر حقيقي.
    \end{obox}
    \begin{obox}{\faExclamationTriangle\quad لا تكتفِ بالنتيجة الأولى}
        النتيجة الأولى في محرك البحث ليست بالضرورة الأفضل:
        قد تكون إعلانًا، أو موقعًا مشهورًا لكنه غير دقيق.
    \end{obox}
\end{frame}

\begin{frame}{معايير تقييم المصدر}
    قبل أن نعتمد على مصدر، نطرح خمسة أسئلة:
    \begin{center}
    \begin{tabularx}{\textwidth}{|>{\centering\arraybackslash}p{2.6cm}|R|}
        \hline
        \ras{المعيار} & \ras{السؤال} \\ \hline
        \textbf{الكاتب}  & من كتب هذا المحتوى؟ هل هو مختص في المجال؟ \\ \hline
        \textbf{الجهة}   & من نشره؟ جامعة، مجلة علمية، أم موقع مجهول؟ \\ \hline
        \textbf{التاريخ} & متى نُشر؟ هل المعلومة ما زالت صالحة؟ \\ \hline
        \textbf{الدقة}   & هل يذكر مراجعه؟ هل نجد المعلومة نفسها في مصدر آخر؟ \\ \hline
        \textbf{الهدف}   & هل هدفه الإعلام، أم الإقناع، أم البيع؟ \\ \hline
    \end{tabularx}
    \end{center}
\end{frame}

\begin{frame}{تلخيص المعلومات}
    \begin{tbox}{\faLightbulb\quad تعريف}
        \textbf{التلخيص:} إعادة كتابة الأفكار الأساسية لنص ما،
        بكلمات أقل وبأسلوبنا الخاص، دون تغيير المعنى.
    \end{tbox}
    \begin{center}
        \flow[1.95cm]{{اقرأ النص\\كاملًا}, {حدد الأفكار\\الأساسية}, {احذف التكرار\\والأمثلة}, {اكتب\\بأسلوبك}, {اذكر\\المصدر}}
    \end{center}
\end{frame}

\begin{frame}{الذكاء الاصطناعي والتلخيص}
    يمكن أن نطلب من الأداة تلخيص نص طويل،
    لكن يجب أن نقارن الملخص بالنص الأصلي:
    \begin{obox}{\faExclamationTriangle\quad أسئلة للمراجعة}
        \begin{itemize}
            \item هل حذفت الأداة فكرة مهمة؟
            \item هل أضافت معلومة غير موجودة في النص؟
        \end{itemize}
    \end{obox}
\end{frame}

\begin{frame}{خلاصة الدرس}
    \begin{tbox}{\faCheckCircle\quad خلاصة الدرس}
        \begin{itemize}
            \item نبدأ البحث بــ\textbf{سؤال محدد}، لا بموضوع عام.
            \item نختار \textbf{كلمات مفتاحية} دقيقة ونضيف مرادفاتها.
            \item نقيّم المصدر: الكاتب، الجهة، التاريخ، الدقة، الهدف.
            \item نلخّص بأسلوبنا، ونذكر \textbf{المصدر} دائمًا.
        \end{itemize}
    \end{tbox}
\end{frame}

\end{document}
